\documentclass[12pt]{article}
\usepackage[utf8]{inputenc}
\usepackage[T1]{fontenc}
\usepackage{hyperref}
\title{Tseku Language Resources for K-12 Curriculum Development}
\author{Compiled by Claude Sonnet 5 (claude-sonnet-5) for Andr\'as Kornai}
\date{}
\begin{document}
\maketitle

\section{Introduction}
Tseku (also spelled Tzuku, ISO 639-3 tsk) is described by Wikipedia as a Tibetic language and, following Tournadre (2013), as part of Khams Tibetan. Speaker figures in the results differ: about 19,000 (Wikipedia, about 2000 data) and 21,000 (Joshua Project, with 14,000 in China and 6,900 in Bhutan). Joshua Project reports no known writing system. I found very little, and no educational materials.

\section{Primary Educational Resources}
None found.

\section{Linguistic and Reference Materials}
\begin{itemize}
\item Tseku language, Wikipedia: \url{https://en.wikipedia.org/wiki/Tseku_language}
\item Glottolog map page for the Tseku language: \url{https://glottolog.org/resource/languoid/id/tsek1238.bigmap.html}
\end{itemize}
A search for the alternate name ``Sd\"erlu'' returned nothing, so I cannot say whether it is a name for this language.

\section{Religious and Community Resources}
\begin{itemize}
\item Joshua Project, Tseku people in China: \url{https://joshuaproject.net/index.php/people_groups/19656/CH}
\item Joshua Project, language tsk (says Bible translation is needed): \url{https://joshuaproject.net/index.php/languages/tsk}
\end{itemize}

\section{Digital and Audio Resources}
None found.

\section{Curriculum Development Context}
The language is undocumented in the sources I could find, and I found no orthography or teaching materials. Any curriculum would start with documentation, working with the community. The identification in the task and the Wikipedia classification should be reconciled before the work is planned.

\end{document}